% Section 1: Introduction. Five paragraphs, following paper/outline.md.
% All result numbers are \Num... macros from paper/numbers.tex.
% The few digits typed below are literature values quoted with their source, not model results.
\section{Introduction}

Phenotyping is widely described as the limiting step in crop breeding \cite{Song2021Highthroughput,Fasoula2020Phenotyping,Kholova2021Pursuit}. A review of crop improvement in the Global South identifies phenotyping as a persistent bottleneck and notes that the path from a program decision to a widely grown variety rarely takes less than ten years, even for beans \cite{Kholova2021Pursuit}. Large-scale breeding phenotyping remains mostly rating by breeders \cite{Walter2015Plant}, and ratings depend on who scores. Rankings of senescence and vigor changed with the staff doing the scoring \cite{ZamanAllah2015Unmanned}, and trained pathologists who scored common bacterial blight in common bean differed from one another \cite{Lasdun2024Participatory}. In a review of a century of visual severity assessment, trained raters typically reached a Lin's concordance of 0.85 to 0.95, whereas inexperienced raters without aids could fall to 0.60 or below \cite{Bock2021Plant}. Conventional phenotyping of common bean is slow and labor-intensive, and imaging with red, green and blue (RGB) cameras has been judged the most cost-effective option for biomass and phenology \cite{Javornik2025Phenotyping}.

Image-based tools and artificial intelligence (AI) tools are validated against manual references and often by run time, and none of the tool papers we read reports a selection-gain or economic endpoint. Studies of plant height in faba bean from unmanned aerial vehicles (UAV) \cite{Mohammadi2024Enhancing}, an on-device maize app \cite{Liu2021Pocketmaize}, image-based scoring of soybean frogeye leaf spot \cite{McDonald2022Automated} and a cloud platform for several crops \cite{Hu2025Openpheno} report agreement statistics and, in some cases, processing speed. Reviews likewise frame progress as accuracy, throughput and computing cost, and a catalog of plant image analysis software listed 172 tools in 2018 \cite{Mochida2018Computer}. Agreement with a manual score need not mean better selection. The manual reference can itself be the weak measurement \cite{McDonald2022Automated}, and two studies tested the downstream consequence directly. In a maize association panel, image-based tassel traits had lower heritability than manual traits yet showed no detectable loss of power in genome-wide association \cite{Gage2018Comparing}. In maize and sorghum, a substantial part of image-extraction error was under genetic control, so more replication does not remove it, and association hits overlapped only partly even when automatic and ground-truth height estimates agreed closely \cite{Zhou2021Identification}. What a more accurate or faster tool is worth to a breeding program therefore remains unanswered.

Breeding economics points to cycle time as the most cost-effective lever for annual gain. The breeder's equation divides genetic gain per cycle by the cycle length, $\Delta G = \sigma_A\, i\, r / L$, where $\sigma_A$ is the additive genetic standard deviation, $i$ the selection intensity, $r$ the selection accuracy and $L$ the years per cycle \cite{Cobb2019Enhancing,Hickey2017Genomic}. Gain rises only with the square root of heritability, and larger candidate populations raise the selection differential in diminishing increments, whereas cycle time is ``the most powerful parameter'' for increasing genetic gain \cite{Cobb2019Enhancing}. For self-pollinated crops of the developing world, shortening the cycle is reported to be more cost-effective than enlarging populations or raising the selection differential \cite{Atlin2022Simple}. In rice, a cost index of gain fell as cycle length fell under rapid generation advance \cite{Bonnecarrere2019Economic}, and shorter recycling schemes raised medium-term gain in stochastic simulation \cite{Seck2024Stochastic}. In Zimbabwean maize under a fixed budget, a restructured pipeline built on doubled haploids raised gain per year and per dollar \cite{Chaingeni2025More}. A dry-bean simulation examined a shortening of the cycle from eight to five years through speed breeding \cite{Lin2023Simulations}, and rapid-cycle genomic selection produced realized yield gain in CIMMYT wheat \cite{Dreisigacker2023Results}. Two limits apply. Under a fixed budget, more cycles per year shrink the population carried through each cycle \cite{Gorjanc2018Optimal}, and rapid-cycling gain plateaus unless the training data continue to be phenotyped \cite{Pook2025Optimization}. Phenotyping throughput is therefore a precondition for short cycles, not an alternative to them.

The only economic model of phenotyping adoption we found is theoretical, and its authors report that accurate information on costs and on the contribution to genetic gain is scarce \cite{Awada2018Adoption}. Programs that deploy field tools state, without quantifying, that reduced error raises gain and shortens cycles \cite{Lasdun2024Participatory}. Mobile-phone imaging tools are being developed for public breeding programs and already used to image on-farm trials in Tanzania \cite{Lasdun2024Participatory}, and a costing study of maize pipelines in Eastern and Southern Africa recommends digital data capture to cut labor cost without valuing it \cite{Das2025Costing}. Stochastic simulations have compared breeding pipelines with and without high-throughput phenotyping by genetic gain and cost-effectiveness \cite{Peixoto2023Simulationbased,Jahufer2021Deterministic}, and a simple simulation showed that higher throughput can compensate for lower measurement accuracy when the extra capacity goes to more environments or larger populations \cite{Lane2021High}. Simulation-optimization of a breeding program under a budget has also been demonstrated \cite{Hassanpour2024Optimization}. None of these studies values a phenotyping tool through variety adoption and discounting, or asks which attribute of a tool a pilot should measure. Ex ante evaluations of breeding investments use adoption paths and net present value (NPV) \cite{Staver2020ExAnte,Grovermann2022Three,Janssen2018Costbenefit}, but such evaluations tend toward optimism: reported rates of return to agricultural research are inflated \cite{Hurley2014Reexamining}, and cost and benefit estimates for other public investments are biased and fat-tailed \cite{Flyvbjerg2021Costbenefit}. This argues for pairing any valuation with sensitivity analysis and a value-of-information framing. To our knowledge, and within the searches described in the Methods, no previous study values AI-based phenotyping through a coupled breeding-simulation and adoption-economics model, or applies value-of-information analysis to this decision.

This study values a generic smartphone computer-vision phenotyping tool, described by its technical attributes rather than by any specific product. A smartphone camera is a proximal sensor, and AI extracts the trait values from its images. In the tool papers we read, such proximal sensing and AI tools are evaluated by agreement with manual measurements, and the model here supplies a decision-relevant endpoint instead: the genetic gain and the economic value a tool with stated properties would add to a breeding program. A stochastic simulation of a four-stage East African common-bean pipeline with a \NumCfgCycleYears-year cycle compares manual and tool phenotyping under one fixed budget across \NumNcellsOK tool profiles. The simulated gain is calibrated to realized progress in common bean and carried through logistic adoption and discounting over a \NumCfgHorizon-year horizon and a grid of economic states. We ask how much each of five technical inputs (measurement-error reduction, whole-plot cost reduction, cycle compression, tool fixed cost and program budget) changes genetic gain, where cycle compression denotes the years by which the tool shortens the breeding cycle; how that gain translates into value across economic states; and which uncertainty a pilot should resolve first. Three meta-analyses of realized gain in common bean, tool-manual agreement, and time and cost savings ground the inputs. The analysis is ex ante and conditional: it reports what a tool with stated properties would be worth in the simulated program, not what a deployed tool has delivered, and cycle compression is an assumed input, not a result. We also propose, and demonstrate on one case, an evaluation chain that runs from tool metrics (agreement, time, cost) through program-level gain under a fixed budget, calibrated gain and adoption-weighted value to attribution and value of information.
